% CCSS 7.RP.A.3 — Simple Interest (FL/TX: financial literacy emphasis)
\section{Simple Interest: Earning and Paying Interest}

\begin{learningGoals}
\begin{itemize}
    \item Use $I = Prt$ to calculate simple interest for savings and loans
    \item Compare financial options using interest calculations
\end{itemize}
\end{learningGoals}

\begin{conceptBox}{The Simple Interest Formula}
\textbf{Simple interest} is money earned on savings or owed on loans.
$$I = Prt$$
\begin{itemize}[leftmargin=*]
    \item $I$ = interest (dollars earned or owed)
    \item $P$ = principal (starting amount)
    \item $r$ = annual rate (as a decimal)
    \item $t$ = time (in years)
\end{itemize}
\textbf{Total amount} = $P + I$
\end{conceptBox}

\begin{workedExample}{Comparing Two Savings Accounts}
You have $\$1{,}000$ to deposit. Bank~A offers $3\%$ simple interest. Bank~B offers $4.5\%$. You plan to save for $2$ years. Which bank earns more?

\textbf{Bank A:} $I = 1{,}000 \times 0.03 \times 2 = \$60$. Total $= \$1{,}060$.

\textbf{Bank B:} $I = 1{,}000 \times 0.045 \times 2 = \$90$. Total $= \$1{,}090$.

\answerHighlight{Bank B earns $\$30$ more. A higher rate means more interest.}
\end{workedExample}

\begin{realWorld}[Savings vs.\ Borrowing]
Interest works in two directions. When you \textbf{save}, the bank pays \textbf{you} interest — your money grows. When you \textbf{borrow}, you pay the bank interest — you owe more than you borrowed. Choosing a higher savings rate or a lower loan rate can save you hundreds of dollars.
\end{realWorld}

\mascotSays{Always convert the percent rate to a decimal first! $4\% = 0.04$, not $4$. Forgetting this gives an answer that is $100$ times too large!}

\begin{practiceBox}{Simple Interest}
\resetProblems

\prob You deposit $\$800$ at $5\%$ for $2$ years. How much interest do you earn?
\answerExplain{$\$80$}{Use $I = Prt$: $I = 800 \times 0.05 \times 2 = \$80$. This is the money the bank pays you for keeping your savings there.}

\prob You borrow $\$1{,}200$ at $6\%$ for $4$ years. How much interest do you owe?
\answerExplain{$\$288$}{$I = 1{,}200 \times 0.06 \times 4 = \$288$. This is the extra cost of borrowing the money.}

\prob Account~A: $\$500$ at $3\%$ for $3$ years. Account~B: $\$500$ at $2\%$ for $5$ years. Which earns more interest?
\answerExplain{Account B ($\$50$ vs.\ $\$45$)}{A: $I = 500 \times 0.03 \times 3 = \$45$. B: $I = 500 \times 0.02 \times 5 = \$50$. Even though B has a lower rate, the longer time makes it earn $\$5$ more.}

\prob A loan of $\$2{,}000$ charges $\$240$ interest over $3$ years. What is the rate?
\answerExplain{$4\%$}{Solve for $r$: $240 = 2{,}000 \times r \times 3$. So $6{,}000r = 240$ and $r = 0.04 = 4\%$.}

\prob You want to earn $\$150$ in interest on $\$1{,}500$ at $2.5\%$. How many years?
\answerExplain{$4$ years}{Solve for $t$: $150 = 1{,}500 \times 0.025 \times t$. So $37.5t = 150$ and $t = 4$ years.}

\end{practiceBox}
